\begin{abstract}
Model merging has emerged as a powerful paradigm for combining multiple task-specific expert models without the prohibitive cost of joint training. While heuristic parameter averaging (e.g., Task Arithmetic) dominates the literature, it often violates the underlying representation geometry, causing severe representation drift. In this work, we present a rigorous theoretical and empirical investigation into the \textbf{limits of representational isotropy on curved manifolds}, using the orthogonal group $\mathrm{O}(d)$ and its associated Lie algebra tangent space $\mathfrak{so}(d)$ as a diagnostic environment. We uncover and formalize two critical phenomena that dictate the success and limits of manifold-level merging. First, we demonstrate that manifold merging is highly sensitive to parameter non-orthogonality; incorporating soft orthogonal regularization during training keeps parameters near the manifold and reduces the Euclidean residual, boosting merged accuracy from $42.07\%$ to $84.55\%$, while attempting to bypass orthogonal pre-training via a post-hoc SVD projection of an unconstrained model is functionally destructive, collapsing merged performance to $15.00\%$. Second, we expose a surprising and profound \textbf{spectral balancing pitfall in Lie algebra spaces}: while spectral balancing (e.g., SAIM) is safe in Euclidean space, performing it in the Lie algebra tangent space $\mathfrak{so}(d)$ (e.g., RIMO) catastrophically scrambles representations, collapsing accuracy to $13.66\%$ on MLP and $18.44\%$ on a Vision Transformer (ViT). We prove via the \textbf{Kernel Distortion Theorem} and \textbf{Spectrum Distortion Theorem} that standard SVD solvers introduce non-symmetric coordinate gauges in the multi-dimensional null space, and that the subsequent skew-symmetric projection step required to restore algebraic closure inevitably distorts the active spectrum. Under the non-linear forward Cayley map, inflating smaller singular values translates to massive, spurious high-dimensional rotations across inactive planes. To overcome this, we propose and validate \textbf{rank-preserving spectral pruning} as a mathematically consistent alternative that keeps inactive dimensions at exactly zero, successfully bypassing the pitfall and achieving robust merged accuracies of $91.49\%$ (MLP) and $88.16\%$ (ViT). Our findings establish rigorous boundaries for Riemannian model merging, warning against naive spectral operations in tangent spaces and highlighting the essential role of native, manifold-respecting models.
\end{abstract}
